\subsection{One hyperfine spectrum, charge profiles and complete gauge records}
\label{sec:hyperfine-charge-gauge-records}

\paragraph{One native hyperfine action.}
A positive current interaction
$h[\tfrac32N_qN_{\bar q}+2\mathbf S_q\cdot\mathbf S_{\bar q}]$
is the sum of quark--antiquark pair triplet squares. It vanishes on every
pure-quark state and shifts a spin-one one-pair meson by $2h$, leaving its
spin-zero partner unchanged. Together with a calibrated coefficient of the
existing local spin square, one declared action gives a neutral vector
energy $0.77526\,\mathrm{GeV}$ and an actual $S=I=3/2$, charge-one baryon
energy $1.232\,\mathrm{GeV}$, while preserving the five specified proton,
neutron, pion, axial and nuclear-binding calibration agreements. These are
seven empirical inputs. The two added mass conventions are Breit--Wigner
numbers, not derived complex resonance poles or widths.
The complete $120$, $64$ and $2520$ baryon, meson and nuclear singlet
spectra, their coloured preparation failures and all magnetic partners are
retained. Exact ground-line bounds test transfer of the native current and
charge records, without identifying states from fitted labels alone.
For the complete local $2^{24}$ matter Fock space, exact SU(3) fusion gives
$\sum_x C_x^2\leq(22/7)\sum_e E_e^2$ on the physical zero-boundary-flux
Gauss carrier. An independent quark/antiquark-count bound also retains the
positive pair and baryon-density terms. Both apply to the same conditional
quadrupole/string-length attachment; neither is a four-dimensional
pure-gauge mass-gap assertion.
The complete source and instruments are specified in \path{model/openwave_common_hyperfine_spectrum}.

\paragraph{Native charge profiles and their independent radius test.}
On the same seven-input strong source, the literal native isospin raising
operator transports the complete neutron ground-ray certificate to the
proton ray. Their mass difference is $0.00129333251$ GeV. The intrinsic
neutron charge is positive at graph radius $1/3$ and negative at $2/3$.
A maximally mixed input on all $120$ physical colour singlets gives a
ground-ray filter probability $1/120$; all $2024$ occupation effects and
their failure records are retained through a compressed positive-operator
measurement. Its nonselective mean source work on that input is zero.

With the shared quadrupole/string length near $0.7559087287$ fm, the
orientation-averaged charge records give a proton root-mean-square radius
near $0.3345681563$ fm and neutron charge second moment near
$-0.03731193543$ fm$^2$. The common length cancels in
$-\langle r_n^2\rangle/\langle r_p^2\rangle$. Its certified interval is
disjoint from the comparison box formed by the CODATA2022 proton radius
and PDG2025 neutron mean-square radius. No single length therefore fits
both radii on this selected source. The graph orientation average is
specified; no relativistic Sachs-current or continuum nucleon-pole
identification follows from these records
(\texttt{model/openwave\_baryon\_form\_factors}).
The complete source and instruments are specified in \path{model/openwave_baryon_form_factors}.

\paragraph{Complete strong-source reduction and a linear potential.}
On the same seven-input action, the complete native baryon singlet carriers
have dimensions $120,364,816,1540$ on two through five cells. Eliminating the
full long-span complement yields the exact Schur action
$H_{\rm eff}(E)=A+V(E-D)^{-1}V^\dagger$, together with its normalization
metric and native current vertex $R_f^\dagger J R_i$.
Rational complement bounds control their variation across a stated energy
window; full finite-time records retain the complement and all failures.
Graph-dependent masses and axial responses are diagnostics, without a
continuum QCD running-coupling claim.
For a pinned quark and diquark at nonadjacent interior sites, all $40$
internal singlet states obey
$H_{\rm pin}(r+1)-H_{\rm pin}(r)=(4\kappa/3)I$.
This exact total-source identity includes local, spin, isospin, hopping and
Gauss-electric terms. Its common conditional length gives a constrained
inter-charge slope $1\,\mathrm{GeV/fm}$, explicitly used as an attachment
target. The result concerns prescribed positions and does not identify the
unpinned ground energy or a four-dimensional pure-gauge mass gap.
The complete source and instruments are specified in \path{model/openwave_strong_step_scaling}.

\paragraph{Calibrated spatial rest bands.}
The complete native three-quark Gauss source admits an infinite-chain
rest-band calculation through five compact cluster states per translation
cell. The omitted relative configurations are retained by a lower spectral
bound, and the exact squared compact-to-complement coupling norm is $5/64$.
Two declared degree-dependent colour and number counterterms vanish on the
entire calibrated two-cell operator. Using the inherited neutron and axial
targets, the complete bulk neutron energy lies in
$[0.939565387,0.939565457]$ GeV and its native axial magnitude lies in
$[1.274757,1.275843]$; the vector matrix element is exactly one. The proton
interval follows by the exact native isospin shift. These are calibrated
one-dimensional source bands with complete spectral complements, without a
three-dimensional scattering-pole or continuum identification. The fixed
coefficient control has a disjoint axial interval
$[1.291676,1.292715]$, so preservation of two-cell fits alone does not imply
spatial refinement.
The complete source and instruments are specified in \path{model/openwave_hadron_refinement_flow}.

\paragraph{Strange families and the complete decay threshold.}
The native colour reduction extends to two quark families with $816$
complete three-quark singlets and $256$ complete quark--antiquark singlets.
Two positive number/pair coefficients calibrate a charged kaon to
$0.493677\,\mathrm{GeV}$ and a Lambda state to $1.115683\,\mathrm{GeV}$.
The full strong action has an explicit charge-conjugation map, which gives
the charged-kaon partner the same energy. This completion includes a
nonlocal total-quark/antiquark-number correction; no range-one locality
is claimed for that term. The Sigma, Xi and neutral-kaon masses are
retained failed comparisons. The native charged-current
construction has a nonzero $\Lambda\to p+\pi^-$ vertex, including every
native five-body output. Its complete two-cell strong final-channel floor
is $1.208108251046\,\mathrm{GeV}$, $92.425251\,\mathrm{MeV}$ above the
Lambda calibration target. An exact variational upper bound on the
initial Lambda energy certifies an actual channel gap greater than
$92.349863\,\mathrm{MeV}$. The full weak norm bound preserves this
positive gap.
Thus a nonzero finite vertex does not supply an energy-conserving strange
decay. The quoted mass fits are strong-source eigenenergies, without an
interacting-pole or lifetime identification.
The complete source and instruments are specified in \path{model/openwave_strange_composites}.

\paragraph{Complete light-nuclear bands.}
The common native action admits a range-three interaction
\[
 B\sum_x P_{N_x=N_{x+1}=N_{x+2}=3}
 \left[(\mathbf S_x+\mathbf S_{x+1}+\mathbf S_{x+2})^2-\frac74\right],
 \qquad B=\frac{865442239271}{68719476736}.
\]
It vanishes identically for total quark number at most six and therefore
preserves the seven two-cell light calibration domains. One additional
empirical input, the triton rest energy, sets $B$. Exact local exterior
characters and colour fusion retain all $572572$ and $176729280$ physical
singlets in three- and four-cell boxes. Their low bands reduce to $64$
and $256$ coordinates with complete complement bounds; rational inertia
certifies every spin and isospin multiplet. The uniform spectral errors
are below $3.868$ and $15.517\,\mathrm{keV}$, respectively.

The three-cell spin/isospin-$1/2$ ground interval is
$[2.8089192030,2.8089230704]\,\mathrm{GeV}$ and contains its calibration
$2.80892113668\,\mathrm{GeV}$. The helion-like interval
$[2.8076258705,2.8076297379]\,\mathrm{GeV}$ and the four-cell
spin/isospin-zero interval
$[3.7391239822,3.7391394985]\,\mathrm{GeV}$ fail the held-out targets
$2.80839161112$ and $3.7273794118\,\mathrm{GeV}$. Complete preparation
failures and higher spectral outputs are retained on the original
maximally mixed Gauss input. The occupancy-dependent number bound proves
full-Fock stability with the common degree and charge-conjugation
corrections. These finite boxes supply neither an asymptotic breakup
threshold nor a physical halo or binding-curve theorem.
The complete source and instruments are specified in \path{model/openwave_light_nuclear_curve}.

\paragraph{Covariant matter and compact Maxwell refinement.}
The declared compact-field matter source uses the actual native charged
spin coordinates and the positive kinetic operator
$\sqrt{m^2+(\boldsymbol\sigma\cdot\boldsymbol\Pi)^2}-m$.
Its Pauli magnetic commutator is retained. Every integer electric flux,
compensating charge reference and native measurement complement remains in
the source. Complete position--flux records have rational tail and finite
arithmetic bounds, with preparation and absorption work included.
A positive Wilson energy lifts all seven extra free lattice corners.
The free matter and quadratic Maxwell symbols share $c^2=1/128$ and have
controlled spatial refinement. The construction preserves excitation number;
no vacuum pair current, signed-Dirac transfer or interacting Lorentz theorem
is supplied. Charged mass ratios and the GeV anchor retain their input ancestry.
The complete source and instruments are specified in \path{model/openwave_dirac_maxwell_refinement}.

\paragraph{Dressed charge profile and current response.}
A positive-energy Pauli source with a compact Maxwell field retains the
magnetic commutator, both spin blocks and every integer electric flux.
The electric and magnetic coefficients are $1/32$ and $1/8$, with a positive
covariant Wilson coefficient $1/16$. Physical probes use the native charge
$Q=-1$ and retain its opposite quantum reference. A complete charge-position
measurement gives $R_x^2=\langle x^2\rangle$, independently of the current
spectrum. On the four-site plaquette the exact elastic form factor is
$F(k)=1-\langle x^2\rangle(1-\cos k)$. Padded rational residuals certify the
infinite-source ground profile. The regulated electron has
$R_x\simeq0.12118$ cell and maximum first-Brillouin-interval suppression
$1-F\simeq0.02937$. The core is therefore below the spatial resolution needed
to demonstrate an order-one onset. The zero-Wilson controls retain the
smaller heavy-family responses.

The physical dipole $D=Qx$ supplies $J=i[H,D]$ and the required contact
$T=-[D,[D,H]]$. The exact static Ward sum is
$\langle T\rangle=2\sum_n(E_n-E_0)|D_{n0}|^2$.
The contact-inclusive imaginary-frequency response is
\[
 \Gamma(\nu)=2\nu^2\left\langle Dg,
 \frac{H-E_0}{(H-E_0)^2+\nu^2}Dg\right\rangle.
\]
Exact padded resolvent residuals and ground-state bounds give absolute
errors below $10^{-9}$ on four complete infinite-source response values.
For the regulated electron these are approximately
$2.96653\times10^{-5}$, $4.42359\times10^{-4}$,
$3.52069\times10^{-3}$ and $6.74284\times10^{-3}$ at
$\nu=1/64,1/16,1/4,1$, respectively, in electron-energy units.
All individual spectral weights and complete probe outcomes are retained.
The source is an intraband construction; these results supply neither
pair-producing vacuum running nor a physical electron-core fit.
The complete source and instruments are specified in \path{model/openwave_dressed_core_response}.

\paragraph{Autonomous collision clocks and their onset control.}
An even native charged-branch transfer, its Gauss reference and neutral
memory define a complete collision source with phase-insensitive pumping
and two-quantum loss. A three-state battery with energies $m,0,3m$ supplies
an energy-conserving dilation. Every event, inactive occupation and Fock
complement is retained. Exact positive-matrix bounds certify a coherence
frequency $m$ and a nonzero finite linewidth. For arrays of four, eight and
sixteen native branches the first-band linewidths are approximately
$0.0333783$, $0.0139327$ and $0.00602053$ in source energy units. The array
frequency is its constituent gap, not its total rest energy.

For a single charged branch the gap is the inherited particle rest energy.
Its complete equatorial records have the factor
$(1-\gamma\delta)^{k/2}\exp(-imk\delta)$. An exact counterfactual shows
that the first probe prepares this conditional phase independently of the
prior pumped occupation: with the pump removed, the oscillation remains
and its linewidth vanishes. Autonomous particle occupation onset is
therefore established, while autonomous internal phase onset is not
established by this instrument. This distinction concerns the actual
observer program; a vanishing ensemble mean by itself does not exclude
oscillatory shared conditional records.
The complete source and instruments are specified in \path{model/openwave_autonomous_clock_channel}.

\paragraph{Fractional charge and the native quotient-Gauss string.}
The complete native representation obeys
$\exp(2\pi iQ)Z_{\mathrm{SU}(3)}=I$ on every Fock sector. On full
$(\mathrm{SU}(3)\times\mathrm{U}(1))/\mathbb Z_3$ tree links the electric
integer $n_e$ and colour triality $t_e$ satisfy $n_e+t_e=0\pmod3$.
The gauge-invariant dual phase $\exp(2\pi i n_e/3)$ therefore jumps at
an actual quark endpoint by $\exp(2\pi iQ)$. The native up and down lifts
are $2/3$ and $-1/3$. This statement follows from the actual native current
and Gauss law, including the quantum reference.

An exact zero-U1-energy dilation preserves the parent's complete strong
Hamiltonian and its calibrated masses. Every off-diagonal source term is
dressed, including charge-one weak exchanges. Complete finite executions
give occupation, flux and triality records with preparation failures and
nonselective probe work retained. This is a dynamical discrete dual-phase
invariant; an electric eigenstate is delocalized in its conjugate link
angle. A continuous classical texture winding and free-quark rest masses
are not inferred from it.
The complete source and instruments are specified in \path{model/openwave_native_centre_string}.

\subsection{Positive energy response and common gravitational records}

An autonomous two-state geometry register couples to both complete native
eight-level factors through their positive ground-referenced energies.
For \(H_g=\omega(I-\sigma_x)/2\), \(Q=q\sigma_z\), and the positive
three-dimensional harmonic profile \(F\), the source
\(H=H_g+e^QK_s+e^{FQ}K_t\) has an explicit square-root ground energy.
Its mixed source/probe energy derivative is strictly negative.
Rational bounds certify attraction for all 49 excited native pairs at
three separations, with both self energies subtracted. Complete autonomous
records conserve combined energy and retain measurement work.
The positive energy coupling is a declared construction; the reciprocal
scalar source's opposite-charge control is not transferred to it.
Its dimensionless Newton coefficient follows from the mode stiffness,
while physical strength and Einstein dynamics require separate attachment.
The complete construction is \path{model/openwave_universal_energy_geometry}.

\paragraph{Joint transverse light and native clocks.}
A complete three-dimensional Maxwell complex has $882$ directed edges and
$540$ transverse coordinates, with Gauss law imposed at every vertex.
The declared positive-frequency field, two native clocks and the quantum
geometry share one autonomous Hamiltonian. Complete finite records give
arrival delay $0.006354\pm0.000509$, transverse displacement
$-0.003216\pm0.000050$, and Fourier-direction change
$-0.004595\pm0.000057$, in source units. The joint native clock phase ratios
are approximately $0.992837$ and $0.992829$. Photon vacuum, no-arrival,
polarization complements and apparatus work are retained. These are finite
constitutive Maxwell records, with no canonical metric quantum-field or
Einstein-continuum transfer asserted. A separate attachment of the electron
energy and strong length fails both the observed Newton strength and light
speed; it supplies no fitted or derived physical $G$.
The complete source and instruments are specified in \path{model/openwave_transverse_gravity}.
